% Gerado de docs/monografia/05-construcao-e-verificacao.md por scripts/monografia_latex.py.
% Não edite aqui: edite o .md e gere de novo.

\chapter{\texorpdfstring{\mbox{CONSTRUÇÃO} \mbox{E} \mbox{VERIFICAÇÃO}}{CONSTRUÇÃO E VERIFICAÇÃO}}\label{cap:5}

Um sistema de ensino faz afirmações em dois planos. No plano do \textbf{conteúdo}, cada tela afirma algo sobre C: que este trecho imprime 9, que o compilador reclama daquela linha, que esta alternativa é a certa. No plano do \textbf{comportamento}, o sistema afirma que trata cada estudante de um jeito: que quem erra com convicção vê a demonstração, que quem não sabe sobe a escada, que a sessão tem teto. Este capítulo descreve como cada plano foi verificado sem depender da opinião de quem escreveu o conteúdo ou o código — e diz, ao fim, o que essa verificação não cobre.

\section{O processo}

Cada parte do sistema foi desenhada por escrito antes de ser construída: um documento de desenho com a intenção, as restrições, as alternativas consideradas e as decisões tomadas, e um plano em tarefas pequenas, cada uma com os testes que a provam. As decisões tomadas durante a construção, quando o plano não respondia, foram registradas com o motivo e o custo de estarem erradas.

O código foi escrito com \textbf{testes antes do código}: o teste de cada comportamento novo foi escrito e visto falhar antes de a implementação existir, e só então a implementação foi feita até o teste passar. Um teste que passa antes da implementação não prova nada — pode estar testando outra coisa —, e por isso vê-lo falhar é parte do método, e não formalidade. Testes existentes nunca foram alterados para passar; quando uma mudança deliberada de comportamento exigiu mudar um teste, a mudança foi declarada no registro do controle de versão, com o motivo.

\section{O compilador como árbitro do conteúdo}

Todo trecho de código das escadas — passos, exemplos resolvidos e demonstrações — passa por um \textbf{conferidor} automático antes de entrar no sistema. O conferidor monta cada trecho dentro de um programa completo — com as bibliotecas padrão de entrada e saída e de alocação, e com as funções do trecho levadas para fora do \texttt{main} —, compila com o gcc 16.2 (\texttt{-{}std=c17 -{}Wall}) e confere o que o texto afirma:

\begin{itemize}
\item nos passos de \textbf{digitar}, roda o programa e compara a saída inteira com a resposta aceita;
\item nos passos de \textbf{completar} e \textbf{escolher} com código, monta cada linha possível e exige que as marcadas como inválidas sejam recusadas pelo compilador e as demais aceitas;
\item nos exemplos resolvidos com saída, roda e compara;
\item nas \textbf{demonstrações}, roda e compara a saída, ou, quando a demonstração é o que o compilador diz, exige que a mensagem afirmada saia, e de uma linha só.
\end{itemize}

As opções de compilador que uma demonstração pode pedir a mais formam uma lista fechada (\texttt{-{}O1}, \texttt{-{}O2}, \texttt{-{}Wextra}, \texttt{-{}Wnull-{}dereference}, \texttt{-{}fanalyzer}, \texttt{-{}Wpointer-{}arith}): o que está ali vira argumento de comando, e cada opção entrou na lista com um teste que mostra a demonstração que ela confirma e que, sem ela, a mesma demonstração não se confirma. Endereços de memória, que mudam a cada máquina e a cada execução, aparecem nas demonstrações como números de exemplo, e o conferidor compara a forma da saída, não o número. Cada processo tem um limite de tempo, para que um programa que não termina não confira nada.

O conferidor roda dentro da suíte de testes, sobre os 13 tópicos: uma escada cujo texto deixou de corresponder ao que o compilador faz derruba a suíte.

As \textbf{cenas} (3.6) passam pela mesma conferência, com um passo a mais. O que uma cena afirma não é uma saída, mas um estado: depois desta linha, \texttt{idade} vale 15 e \texttt{p} aponta para \texttt{idade}. O conferidor monta, então, um segundo programa — o código da revisão com, em cada ponto por onde a cena passa, as afirmações dela escritas em C — e o executa; a afirmação que falha sai com a linha e com o que a cena dizia. Três cuidados tornam isso possível dentro de funções e no monte. Cada personagem tem o seu endereço guardado assim que nasce, e toda afirmação passa por esse endereço, e não pelo nome: ela vale dentro da função e fora dela, e também para o bloco do \texttt{malloc}, que não tem nome. A mesma linha pode rodar mais de uma vez — a função chamada duas vezes —, e cada passagem é conferida com o seu próprio estado: a cena que dissesse que a variável local de \texttt{conta} continua valendo 2 na segunda chamada seria acusada, e a primeira chamada, que está certa, não. E, no fim do programa, o conferidor confere quantas vezes passou por cada ponto, o que acusa a cena que conta outra história de chamadas. Antes do compilador, a própria estrutura da cena já recusa o que o tópico ensina a não fazer: ler uma caixa sem valor, seguir o bilhete de um \texttt{NULL} ou de um bloco devolvido, citar uma caixa que já sumiu.

As \textbf{questões do banco} passam por uma conferência mais estreita, que também roda na suíte. Nas 321 questões que têm código, ela procura string partida entre duas linhas — uma linha com aspas ímpares, que na tela é o código cortado ao meio. Nas questões de prever a saída que imprimem e cuja alternativa certa começa por um valor (``30: só o último valor escrito''), ela monta o programa, compila, roda e exige aquele valor; a questão que não bate tem de estar fora da trilha. O resto do banco não se confere comparando texto: a alternativa certa que diz ``não dá para saber'', e as questões de achar o erro e de comparar casos, que muitas vezes são fragmentos que não formam programa. Para elas fica a revisão mecânica da forma — toda alternativa errada aponta uma concepção do catálogo, a certa não é visivelmente a mais longa, o enunciado não afirma a severidade de uma mensagem que o compilador pode rebaixar a aviso, e as questões servidas não usam construções ainda não ensinadas (5.3) — e a leitura.

Essa conferência do banco foi escrita depois que a revisão deste texto notou a falta dela. Na primeira rodada, achou de novo a questão do kc\_12 que já estava fora da trilha (depois reescrita), e achou oito questões do kc\_02 com o código partido (6.4).

\section{A régua do texto, em teste}

As regras de escrita da escada (3.4) que podem ser decididas por programa foram transformadas em testes, que rodam sobre todos os tópicos:

\begin{itemize}
\item de três a sete passos por premissa, com os guiados antes dos de prática, e pelo menos dois de prática;
\item no máximo duas frases de até 30 palavras por passo guiado, e nenhuma nos de prática;
\item duas versões diferentes por passo, e todo exercício com erro previsto e retorno de pelo menos cinco palavras;
\item o verbo ``reclama'' para o que o compilador diz, nunca ``não compila'' nem ``recusa'', porque a severidade de uma mensagem muda entre versões do compilador;
\item nenhum passo de digitar cuja resposta seja um endereço;
\item demonstrações de até nove linhas, sem pergunta embutida, uma por concepção errônea do tópico, com a previsão em posições variadas;
\item a concepção errônea dita, numa alternativa marcada, só em passo guiado da premissa que ela nega;
\item todo primeiro passo de prática com um nome curto, que é o que o fechamento mostra;
\item o treino depois da demonstração não repete o trecho que a demonstração acabou de mostrar;
\item as questões servidas pela trilha não usam construções que nenhum tópico ensinou até ali, e cada premissa abre com questão de código.
\end{itemize}

Para cada tópico, há ainda testes do que ele apresenta e onde: em que passo guiado cada construção nova aparece pela primeira vez (\texttt{static}, a conversão \texttt{(char *)}, o ponteiro para ponteiro), e quais questões do banco ficam fora da trilha, com o motivo de cada uma.

\section{Os estudantes simulados}

O comportamento foi verificado com \textbf{dez estudantes simulados}, cada um um perfil de resposta definido por programa: o que sabe tudo; o inseguro, que acerta sem certeza; o que não sabe e aprende com os passos; o que não sabe e erra todos os passos; o que sustenta a concepção errônea de cada premissa e muda de ideia depois da demonstração; o que sustenta e não muda; o que escorrega, errando com a razão certa; o que chuta com certeza; o que tropeça na prática; e o que abre a explicação antes de responder.

Cada perfil atravessa cada tópico do começo ao fim, nos dois ritmos, pelo mesmo motor que atende as pessoas. Os testes conferem, nessas travessias, as propriedades que o desenho afirma: toda premissa termina com um dos cinco resultados; nenhuma premissa recebe mais de um trabalho; a sessão tem teto; quem sustenta uma concepção errônea vê a demonstração dela; a sessão inteira pode ser reconstruída a partir do registro; e perfis diferentes terminam com resultados diferentes. A mesma travessia produz a \textbf{medida} de cada tópico — telas e toques por perfil —, que o Capítulo 6 apresenta.

\section{A interface}

Também a interface foi verificada por teste onde isso foi possível: a razão de contraste de cada par de cores de texto e fundo, nos dois temas, calculada pelas definições de luminância relativa e de razão de contraste da WCAG 2.1 \cite{w3c2025}; que o realce de código devolve exatamente o texto original, em todos os trechos de código das escadas; que nenhuma página carrega nada de outro endereço; que o tema é escolhido antes de a página aparecer, inclusive em navegador sem armazenamento local; e que a saída do programa nunca recebe cores de código. O teste de contraste usa o limiar de linearização da versão 2.0 da norma, 0,03928, que a versão atual trocou por 0,04045, dizendo que a troca não tem efeito prático: em cores de 8 bits nenhum valor cai entre os dois, e outro teste confere isso. Das cenas, o que se testa por programa são os estados que o desenhista calcula a cada passo — que a caixa da função some quando ela termina, que o bloco devolvido fica apagado, que a frase de cada passo existe para o leitor de tela. O resto — o desenho de cada tela e de cada personagem, o comportamento em celular (360 px de largura) e em tela larga — foi conferido no navegador, tela a tela, nos dois temas.

\section{O ensaio}

Cada tópico foi atravessado também por sessões de ensaio, num servidor separado do que está no ar e gravando numa pasta separada da de coleta, para que nenhuma sessão simulada se misturasse a dados de pessoas. As sessões foram levadas pelo protocolo HTTP, pelos mesmos perfis simulados, até cada tipo de tela, e conferidas na página.

\section{A leitura das perguntas}

O conferidor diz se uma afirmação é verdadeira, e não se a pergunta é boa. Quando o autor leu o primeiro passo do kc\_02 no sistema, achou a pergunta confusa: depois de duas frases que diziam o que é o endereço de uma variável, ela perguntava ``o que é o endereço de \texttt{idade}?'', e a alternativa certa era a frase de cima copiada — responder era achar palavras, e quem sabia desconfiava de pegadinha. O passo foi reescrito com um endereço de exemplo (``\texttt{idade} ficou no lugar de número 1000 e guarda 15: qual é o endereço de \texttt{idade}?''), e o mesmo defeito foi procurado em todos os tópicos.

O critério escrito foi este: há defeito quando a frase do passo diz a conclusão sobre o próprio caso da pergunta — o mesmo código, os mesmos nomes —, ou quando a última frase é a alternativa certa. Não há quando a frase dá uma regra geral e a pergunta obriga a reconhecer, no código, qual coisa é qual antes de aplicá-la; nem no passo que mostra uma construção nova e pede para imitá-la uma vez, que é assim de propósito. Um programa marcou os passos em que as palavras da alternativa certa já estavam nas frases, e os 244 passos guiados foram lidos, os marcados e os demais. O resultado está em 6.4. Na reescrita, os testes do que cada tópico apresenta (5.3) serviram de rede: cinco frases novas tinham perdido a construção que aquele passo deve apresentar, e os testes as acusaram.

\section{O que esta verificação não cobre}

O compilador garante que o que a tela afirma sobre um programa é verdade. Ele não garante que a ideia ensinada é a ideia certa, que a escada está na ordem mais fácil para um iniciante, que a analogia ajuda, ou que a concepção errônea do catálogo é a que um estudante de fato tem. Os testes garantem que o sistema se comporta como o desenho diz; não garantem que o desenho ensina. A primeira lacuna é preenchida pela revisão do conteúdo por quem conhece o domínio — o conteúdo das escadas foi escrito com assistência de um modelo de linguagem e ainda aguarda a revisão do autor \textbf{[atualizar quando a revisão for feita]} —; a segunda, só por um estudo com estudantes, que é o primeiro dos trabalhos futuros (Capítulo 7).
